%  LaTeX support: latex@mdpi.com 
%  For support, please attach all files needed for compiling as well as the log file, and specify your operating system, LaTeX version, and LaTeX editor.

%=================================================================
\documentclass[jimaging,article,submit,pdftex,moreauthors]{Definitions/mdpi} 

\graphicspath{{figures/}} % path to folder with figures
\usepackage{caption}
\usepackage{subcaption}
%\usepackage[export]{adjustbox}
\usepackage{listings}
\usepackage{xcolor}
\usepackage{graphicx}
\usepackage{gensymb} % for degree symbol

\definecolor{codegreen}{rgb}{0,0.6,0}
\definecolor{codegray}{rgb}{0.5,0.5,0.5}
\definecolor{codepurple}{rgb}{0.58,0,0.82}
%\definecolor{backcolour}{rgb}{0.95,0.95,0.92}
\definecolor{backcolour}{RGB}{255,250,240}

\lstdefinestyle{mystyle}{
    backgroundcolor=\color{backcolour},   
    commentstyle=\color{codegreen},
    keywordstyle=\color{magenta},
    numberstyle=\tiny\color{codegray},
    stringstyle=\color{codepurple},
    basicstyle=\ttfamily\footnotesize,
    breakatwhitespace=false,         
    breaklines=true,                 
    captionpos=t, 
    keepspaces=true,                 
    numbers=left,                    
    numbersep=5pt,                  
    showspaces=false,                
    showstringspaces=false,
    showtabs=false,                  
    tabsize=2
    }

\lstset{style=mystyle}
\captionsetup[lstlisting]{position=top, margin=0cm,labelfont={bf, small, stretch=1.17}, labelsep=colon, textfont={small, stretch=1.17}, aboveskip=6pt, singlelinecheck=off, justification=justified}
	
\usepackage{soul} % highlighting text

%--------------------
% Class Options:
%--------------------
%----------
% journal
%----------

%---------
% article
%---------

%=================================================================
% MDPI internal commands
\firstpage{1} 
\makeatletter 
\setcounter{page}{\@firstpage} 
\makeatother
\pubvolume{1}
\issuenum{1}
\articlenumber{0}
\pubyear{2022}
\copyrightyear{2022}
%\externaleditor{Academic Editor: Firstname Lastname}
\datereceived{} 
%\daterevised{} % Only for the journal Acoustics
\dateaccepted{} 
\datepublished{} 
%\datecorrected{} % Corrected papers include a "Corrected: XXX" date in the original paper.
%\dateretracted{} % Corrected papers include a "Retracted: XXX" date in the original paper.
\hreflink{https://doi.org/} % If needed use \linebreak
%\doinum{}
%------------------------------------------------------------------
% The following line should be uncommented if the LaTeX file is uploaded to arXiv.org
%\pdfoutput=1

%=================================================================
% Add packages and commands here. The following packages are loaded in our class file: fontenc, inputenc, calc, indentfirst, fancyhdr, graphicx, epstopdf, lastpage, ifthen, lineno, float, amsmath, setspace, enumitem, mathpazo, booktabs, titlesec, etoolbox, tabto, xcolor, soul, multirow, microtype, tikz, totcount, changepage, attrib, upgreek, cleveref, amsthm, hyphenat, natbib, hyperref, footmisc, url, geometry, newfloat, caption

%=================================================================
% Full title of the paper (Capitalized)
\Title{Satellite Image Processing by Python and R for Terrain Analysis of Côte d'Ivoire, West Africa}

% MDPI internal command: Title for citation in the left column
\TitleCitation{Satellite Image Processing by Python and R for Terrain Analysis of Côte d'Ivoire, West Africa}

% Author Orchid ID: enter ID or remove command
\newcommand{\orcidauthorA}{0000-0002-5759-1089} % Add \orcidA{} behind the author's name
\newcommand{\orcidauthorB}{0000-0002-6461-1551} % Add \orcidB{} behind the author's name

% Authors, for the paper (add full first names)
\Author{Polina Lemenkova $^{1}$\orcidA{} and Olivier Debeir $^{1}\orcidB{}$}
% \dagger,\ddagger
%\longauthorlist{yes}

% MDPI internal command: Authors, for metadata in PDF
\AuthorNames{Polina Lemenkova $^{1}$ and Olivier Debeir $^{1}$}

% MDPI internal command: Authors, for citation in the left column
\AuthorCitation{Lemenkova, P.; Debeir, O.}
% If this is a Chicago style journal: Lastname, Firstname, Firstname Lastname, and Firstname Lastname.

% Affiliations / Addresses (Add [1] after \address if there is only one affiliation.)
\address{%
$^{1}$ \quad Université Libre de Bruxelles, École polytechnique de Bruxelles (EPB, Brussels Faculty of Engineering), Laboratory of Image Synthesis and Analysis, Building L, Campus de Solbosch, ULB -- LISA CP165/57, Avenue Franklin D. Roosevelt 50, B-1050, Brussels, Belgium}

% Contact information of the corresponding author
\corres{Correspondence: polina.lemenkova@ulb.be; Tel.: +32471860459 (P.L.)}

% Abstract (Do not insert blank lines, i.e. \\) 
\abstract{A single paragraph of about 200 words maximum.}

% Keywords
\keyword{Image processing, remote sensing; programming; Python; satellite image; Digital Terrain Model; cartography; mapping; spatial analysis; environment} 

% The fields PACS, MSC, and JEL may be left empty or commented out if not applicable
%\PACS{J0101}
%\MSC{}
%\JEL{}

%%%%%%%%%%%%%%%%%%%%%%%%%%%%%%%%%%%%%%%%%%
\begin{document}

\section{Introduction}
\newpage
\section{Study Area}

\begin{figure}[H]\centering
	\includegraphics[width=12.0 cm]{Fig_01}
	\caption{Topographic map of Côte d'Ivoire. Mapping software: Generic Mapping Tools (GMT) scripting toolset. Data source: GEBCO/SRTM. Cartography: authors.
	\label{fig1}}
\end{figure}  

%%%%%%%%%%%%%%%%%%%%%%%%%%%%%%%%%%%%%%%%%%
\section{Materials and Methods}

\subsection{Data}
The satellite image used for data processing is Landsat 9 OLI/TIRS C2 L1 scene covering the central part of Côte d'Ivoire, \emph{LC09\_L1TP\_197055\_20220111\_20220122\_02\_T1}; Scene Identifier: LC91970552022011LGN01; Land Cloud Coverage: 1.29; Date acquired: 2022/01/11. The data used for terrain analysis is SRTM Void Filled DEM, from the SRTM3 (3 arc-second resolution, 90 m), NASA/NGA. Entity ID: SRTM3N07W006V2; tile coordinates: 5°00'00"-6°00'00"W, 7°00'00"-8°00'00"N. Aquisition: 2000-02-11. The data were downloaded from the USGS GloVis, which is a satellite image repository with public open access to remote sensing data: Figure \ref{fig0}.

\begin{figure}
     \centering
     \begin{subfigure}[b]{0.60\textwidth}
         \centering
         \includegraphics[width=12cm]{Printscreen_3}
         \caption{Collecting Landsat 9 OLI/TIRS C2 L1.}
         \label{fig:00a}
     \end{subfigure} \vfill \vspace{2mm}
     \begin{subfigure}[b]{0.60\textwidth}
         \centering
         \includegraphics[width=12cm]{Printscreen_2}
         \caption{Collecting SRTM Void Filled DEM.}
         \label{fig:00a}
     \end{subfigure}
        \caption{Data capture from the USGS Global Visualization Viewer (GloVis) satellite image repository to accessing remote sensing data: Landsat 9 OLI/TIRS C2 L1 and SRTM DEM for terrain analysis}
        \label{fig0}
\end{figure}

\newpage
\subsection{Methodology}

\begin{lstlisting}[language=bash, caption=Unix shell script by GDAL and echo-grep utilities to inspect the geometry of Landsat TIFs]
for f in *tif ; do  echo $f; gdalinfo $f | grep Size; done
\end{lstlisting}

\begin{figure}
     \centering
     \begin{subfigure}[b]{0.45\textwidth}
         \centering
         \includegraphics[width=7cm]{Fig_00a}
         \caption{Using grep and echo Unix utilities with 'gdalinfo' module of GDAL to scan metadata}
         \label{fig:00a}
     \end{subfigure} \hfill \hspace{2mm}
     \begin{subfigure}[b]{0.45\textwidth}
         \centering
         \includegraphics[width=7cm]{Fig_00b}
         \caption{Creating SpatRaster objects in R and check spatial data: projection, datum, ellipsoid, etc.}
         \label{fig:00a}
     \end{subfigure}
        \caption{Inspecting the geometry of Landsat bands (a) which shows 30 pixels resolution for all the bands except for Band 8 (panchromatic) with 15-pixels resolution. Checking image properties in R package 'terra' (b) for single Landsat-8 OLI/TIRS C1 bands (1:11)}
        \label{fig0}
\end{figure}

\begin{lstlisting}[language=python, caption=Python code (1) used for plotting DEM in Figure 2.]
python3.10
# Set home directory
import os
os.chdir('/Users/polinalemenkova/Documents/Python/Image_Processing')
os. getcwd()
# Set data source
dtm = "n07_w006_3arc_v2.tif"
# Load libraries
from glob import glob
import earthpy as et
import earthpy.spatial as es
import earthpy.plot as ep
import rasterio as rio
import matplotlib.pyplot as plt
import numpy as np
from matplotlib.colors import ListedColormap
import plotly.graph_objects as go
np.seterr(divide='ignore', invalid='ignore')
# Open DEM by Rasterio which reads GeoTIFF format and stores gridded raster datasets of digital terrain models (DTMs) and satellite images
with rio.open(dtm) as src:
    elevation = src.read(1)
# Plot the data
fig, ax = plt.subplots(figsize=(10, 10))
ep.plot_bands(
    elevation,
    ax=ax,
    cmap="terrain",
    title="Ditigal Terrain Model (DTM) of Kossou Lake, C\^ote d'Ivoire. \n Data source: SRTM3 (3 arc-second resolution, 90 m), NASA/NGA. \nEntity ID: SRTM3N07W006V2. Aquisition date: 2000-02-11. cpt: 'terrain'",
    figsize=(10, 6),
)
plt.show()
# Save figure with 300 dpi resolution
plt.savefig('Figure_2_DEM_Kossou_1609.jpg', dpi=300)
\end{lstlisting}

\begin{lstlisting}[language=python, caption=Python code (2) used for plotting hillshade in Figure 3.]
python3.10
# Set home directory
import os
os.chdir('/Users/polinalemenkova/Documents/Python/Image_Processing')
os. getcwd()
# Set data source
dtm = "n07_w006_3arc_v2.tif"
# Load libraries
from glob import glob
import earthpy as et
import earthpy.spatial as es
import earthpy.plot as ep
import rasterio as rio
import matplotlib.pyplot as plt
import numpy as np
from matplotlib.colors import ListedColormap
import plotly.graph_objects as go
np.seterr(divide='ignore', invalid='ignore')
# Generate hillshade from DTM
hillshade = es.hillshade(elevation)
# Plot the data
fig, ax = plt.subplots(figsize=(10, 10))
ep.plot_bands(
    hillshade,
    ax=ax,
    cbar=True,
    cmap="cividis",
    title="Hillshade from DTM of Kossou Lake, C\^ote d'Ivoire. \n Data source: SRTM3 (3 arc-second resolution, 90 m), NASA/NGA. \nEntity ID: SRTM3N07W006V2. Aquisition date: 2000-02-11. cmap=cividis",
    figsize=(10, 10),
)
plt.show()
plt.savefig('Figure_3_Hillshade_Kossou_1609.png', dpi=300)
\end{lstlisting}

\begin{lstlisting}[language=python, caption=Python code (3) and (4) used for plotting hillshade with changed azimuth and sun angle values for Figures 4 and 5.]
python3.10
# Set home directory
import os
os.chdir('/Users/polinalemenkova/Documents/Python/Image_Processing')
os. getcwd()
# Set data source
dtm = "n07_w006_3arc_v2.tif"
# Load libraries
from glob import glob
import earthpy as et
import earthpy.spatial as es
import earthpy.plot as ep
import rasterio as rio
import matplotlib.pyplot as plt
import numpy as np
from matplotlib.colors import ListedColormap
import plotly.graph_objects as go
np.seterr(divide='ignore', invalid='ignore')
# Open DEM by Rasterio which reads GeoTIFF format and stores gridded raster datasets of digital terrain models (DTMs) and satellite images
with rio.open(dtm) as src:
    elevation = src.read(1)
# Changing the azimuth of the sun of the hillshade layer to 210 degrees
hillshade_azimuth_230 = es.hillshade(elevation, azimuth=230)
# Plotting the hillshade layer with the modified azimuth
fig, ax = plt.subplots(figsize=(10, 10))
ep.plot_bands(
    hillshade_azimuth_230,
    ax=ax,
    cbar=True,
    cmap="plasma",
    title="Hillshade from DTM of Kossou Lake, C\^ote d'Ivoire. Azimuth of sun is adjusted to 230 \N{DEGREE SIGN}",
    figsize=(10, 10),
)
plt.show()
plt.savefig('Figure_4_Azimuth_hillshade_Kossou.jpg', dpi=300)
# Changing the angle altitude of sun
hillshade_angle_10 = es.hillshade(elevation, altitude=10)
# Plotting the hillshade layer with the modified angle altitude
fig, ax = plt.subplots(figsize=(10, 10))
ep.plot_bands(
    hillshade_angle_10,
    ax=ax,
    cbar=True,
    cmap="magma",
    title="Hillshade from DTM of Kossou Lake, C\^ote d'Ivoire. Angle altitude is defined as 10 \N{DEGREE SIGN}",
    figsize=(10, 10),
)
plt.show()
plt.savefig('Figure_5_Angle_hillshade_Kossou.jpg', dpi=300)
\end{lstlisting}

\begin{lstlisting}[language=r, caption=R code used for plotting color composites for bands of the satellite image]
library(raster)
blue = raster("LC09_L1TP_197055_20220111_20220122_02_T1_B2.TIF")
green = raster("LC09_L1TP_197055_20220111_20220122_02_T1_B3.TIF")
red = raster("LC09_L1TP_197055_20220111_20220122_02_T1_B4.TIF")
rgb = stack(red, green, blue)
plotRGB(rgb, scale=65535)
plotRGB(rgb, stretch="lin")
plotRGB(rgb, stretch="hist")
# False color composites
# Color Infrared: colors near-infrared (band 5) as red, red (band 4) as green, and green (band 3) as blue
green = raster("LC09_L1TP_197055_20220111_20220122_02_T1_B3.TIF")
red = raster("LC09_L1TP_197055_20220111_20220122_02_T1_B4.TIF")
near.infrared = raster("LC09_L1TP_197055_20220111_20220122_02_T1_B5.TIF")
rgb = stack(near.infrared, red, green)
# Create an empty raster() with the top left (ymx and xmn) and bottom right (ymn and xmx) lat/long coordinates.
#boundary = raster(ymx=6.91, xmn=-8.27, ymn=6.18, xmx=-4.83)
#boundary = projectExtent(boundary, rgb@crs)
#rgb = crop(rgb, boundary)
plotRGB(rgb, scale=65535, stretch="lin")
#
# Landsat 8 OLI colour composite image, RGB=753
green = raster("LC09_L1TP_197055_20220111_20220122_02_T1_B7.TIF")
red = raster("LC09_L1TP_197055_20220111_20220122_02_T1_B5.TIF")
near.infrared = raster("LC09_L1TP_197055_20220111_20220122_02_T1_B3.TIF")
rgb = stack(near.infrared, red, green)
plotRGB(rgb, scale=65535, stretch="hist")
# Landsat 8 OLI colour composite image, RGB=765.
red = raster("LC09_L1TP_197055_20220111_20220122_02_T1_B7.TIF")
green = raster("LC09_L1TP_197055_20220111_20220122_02_T1_B6.TIF")
blue = raster("LC09_L1TP_197055_20220111_20220122_02_T1_B5.TIF")
rgb = stack(red, green, blue)
plotRGB(rgb, scale=65535, stretch="hist")
# Landsat 8 OLI colour composite image, RGB=742.
red = raster("LC09_L1TP_197055_20220111_20220122_02_T1_B7.TIF")
green = raster("LC09_L1TP_197055_20220111_20220122_02_T1_B4.TIF")
blue = raster("LC09_L1TP_197055_20220111_20220122_02_T1_B2.TIF")
rgb = stack(red, green, blue)
plotRGB(rgb, scale=65535, stretch="hist")
# Landsat 8 OLI colour composite image, RGB=573.
red = raster("LC09_L1TP_197055_20220111_20220122_02_T1_B5.TIF")
green = raster("LC09_L1TP_197055_20220111_20220122_02_T1_B7.TIF")
blue = raster("LC09_L1TP_197055_20220111_20220122_02_T1_B3.TIF")
rgb = stack(red, green, blue)
plotRGB(rgb, scale=65535, stretch="hist")
# Landsat 8 OLI colour composite image, Color Infrared: RGB=543
green = raster("LC09_L1TP_197055_20220111_20220122_02_T1_B5.TIF")
red = raster("LC09_L1TP_197055_20220111_20220122_02_T1_B4.TIF")
near.infrared = raster("LC09_L1TP_197055_20220111_20220122_02_T1_B3.TIF")
rgb = stack(near.infrared, red, green)
plotRGB(rgb, scale=65535, stretch="hist")
# Landsat 8 OLI colour composite image, Color Infrared: RGB=654
green = raster("LC09_L1TP_197055_20220111_20220122_02_T1_B6.TIF")
red = raster("LC09_L1TP_197055_20220111_20220122_02_T1_B5.TIF")
near.infrared = raster("LC09_L1TP_197055_20220111_20220122_02_T1_B4.TIF")
rgb = stack(near.infrared, red, green)
plotRGB(rgb, scale=65535, stretch="hist")
\end{lstlisting}

%%%%%%%%%%%%%%%%%%%%%%%%%%%%%%%%%%%%%%%%%%
\newpage
\section{Results}
%--------
\subsection{Color composites}
In Figure \ref{fig7} Bands 3-4-5 were taken in inverse way where near-infrared band (B5) was representing red, red (B4) as green, and green (B3) as blue.

\begin{figure}
\makebox[0.90\linewidth][c]{%
	\begin{subfigure}[b]{.6\textwidth}
		\centering
			\includegraphics[width=0.95\textwidth]{Fig_07a}
		\caption{Bands 2-3-4, stretch improvement: 'linear'}
	\end{subfigure}%
	\begin{subfigure}[b]{.6\textwidth}
		\centering
		\includegraphics[width=0.95\textwidth]{Fig_07b}
		\caption{Bands 2-3-4, stretch improvement: 'histogram'.}
	\end{subfigure}%
}\\
\vfill \vspace{1mm}
\makebox[0.90\linewidth][c]{%
	\begin{subfigure}[b]{.6\textwidth}
		\centering
			\includegraphics[width=0.95\textwidth]{Fig_08a}
		\caption{Bands 3-4-5.}
	\end{subfigure}%
	\begin{subfigure}[b]{.6\textwidth}
		\centering
		\includegraphics[width=0.95\textwidth]{Fig_08b}
		\caption{Bands 7-6-5}
	\end{subfigure}%
}
\caption{Natural color composites: (a) and (b). False color composites: (c) and (d). Landsat 9 OLI image of Kossou Lake region and Yamoussoukro, Côte d'Ivoire. Mapping: RStudio. Source: authors.}\label{fig7}
\end{figure}

\begin{figure}
\makebox[0.90\linewidth][c]{%
	\begin{subfigure}[b]{.6\textwidth}
		\centering
			\includegraphics[width=0.95\textwidth]{Fig_09a}
		\caption{Bands 5-7-3.}
	\end{subfigure}%
	\begin{subfigure}[b]{.6\textwidth}
		\centering
			\includegraphics[width=0.95\textwidth]{Fig_09b}
		\caption{Bands 7-4-2}
	\end{subfigure}%
}\\
 \vfill \vspace{1mm}
\makebox[0.90\linewidth][c]{%
	\begin{subfigure}[b]{.6\textwidth}
		\centering
			\includegraphics[width=0.95\textwidth]{Fig_10a}
		\caption{Bands 5-4-3 (Color Infrared).}
	\end{subfigure}%
	\begin{subfigure}[b]{.6\textwidth}
		\centering
			\includegraphics[width=0.95\textwidth]{Fig_10b}
		\caption{Bands 6-5-4}
	\end{subfigure}%
}
\caption{False color composite visualization of Landsat 8 OLI image of Kossou Lake surroundings, Côte d'Ivoire. Mapping: RStudio. Source: authors.}
\end{figure}

\newpage
\subsection{Terrain Analysis}

\begin{figure}
\makebox[0.90\linewidth][c]{%
	\begin{subfigure}[b]{.6\textwidth}
		\centering
			\includegraphics[width=0.95\textwidth]{Fig_02}
		\caption{DTM of Kossou Lake, Côte d'Ivoire.}
	\end{subfigure}%
	\begin{subfigure}[b]{.6\textwidth}
		\centering
			\includegraphics[width=0.95\textwidth]{Fig_03}
		\caption{Hillshade from DTM of Kossou Lake, Côte d'Ivoire.}
	\end{subfigure}%
}\\
 \vfill \vspace{2mm}
\makebox[0.90\linewidth][c]{%
	\begin{subfigure}[b]{.6\textwidth}
		\centering
			\includegraphics[width=0.95\textwidth]{Fig_04}
		\caption{Azimuth of sun: 230\textdegree , hillshade from DTM.}
	\end{subfigure}%
	\begin{subfigure}[b]{.6\textwidth}
		\centering
			\includegraphics[width=0.95\textwidth]{Fig_05}
		\caption{Sun angle altitude: 10\textdegree , hillshade from DTM.}
	\end{subfigure}%
}
\caption{Terrain modelling for Kossou Lake area, Côte d'Ivoire. Mapping: Python. Source: authors.}
\end{figure}

%\newpage
\begin{figure}[H]\centering
	\includegraphics[width=8.0 cm]{Fig_06}
	\caption{SRTM3 Void Filled DEM (NASA/NGA) overlayed on top of a hillshade in Kossou Lake, Côte d'Ivoire, West Africa. cmaps: 'turbo' and 'gist\_yarg'. Mapping: Python. Source: authors.
	\label{fig6}}
\end{figure} 


%%%%%%%%%%%%%%%%%%%%%%%%%%%%%%%%%%%%%%%%%%
\newpage
\section{Discussion}


%%%%%%%%%%%%%%%%%%%%%%%%%%%%%%%%%%%%%%%%%%
\section{Conclusions}

%%%%%%%%%%%%%%%%%%%%%%%%%%%%%%%%%%%%%%%%%%

\funding{This research received no external funding.}

\institutionalreview{Not applicable}

\informedconsent{Not applicable}

\dataavailability{Not applicable} 

\acknowledgments{We thank the two anonymous reviewers for comments and reading of this manuscript}

\conflictsofinterest{The authors declare no conflict of interest} 

\abbreviations{Abbreviations}{
The following abbreviations are used in this manuscript:\\

\noindent 
\begin{tabular}{@{}ll}
CPT & Colour Palette Table \\
DCW & Digital Chart of the World \\
DTM & Ditigal Terrain Model (DTM) \\
GEBCO & General Bathymetric Chart of the Oceans \\
GDAL & Geospatial Data Abstraction Library \\
GIS & Geographic Information System \\
GloVis & Global Visualization Viewer \\
GMT & Generic Mapping Tools \\
Landsat 8-9 OLI/TIRS & Landsat 8-9 Operational Land Imager and Thermal Infrared \\
NASA & National Aeronautics and Space Administration \\
NGA & National Geospatial-Intelligence Agency \\
SRTM & Shuttle Radar Topography Mission \\
USGS & United States Geological Survey \\
\end{tabular}
}

\reftitle{References}

% Please provide either the correct journal abbreviation (e.g. according to the “List of Title Word Abbreviations” http://www.issn.org/services/online-services/access-to-the-ltwa/) or the full name of the journal.
% Citations and References in Supplementary files are permitted provided that they also appear in the reference list here. 

%=====================================
% References, variant A: external bibliography
%=====================================
\bibliography{IvoryCoast.bib}
\nocite{*}

%%%%%%%%%%%%%%%%%%%%%%%%%%%%%%%%%%%%%%%%%%
%\end{adjustwidth}
\end{document}

